\section{Preguntas Metodológicas}

\subsection{¿Cómo se abordará el problema?}
\textbf{Descripción del enfoque adoptado:} Se realizará una revisión documental mediante la recopilación de información proveniente de artículos científicos, bases de datos especializadas y registros históricos. El análisis se centrará en los métodos y técnicas de neutralización, contención y tratamiento de residuos peligrosos tras accidentes nucleares, especialmente aquellos clasificados como nivel 7 en la escala INES.

\subsection{¿Qué tipo de datos se requiere?}
Se requieren datos como:

\begin{table}[!t]
\renewcommand{\arraystretch}{1.2}
\caption{Variables del modelo considerando mitigación y gestión}
\label{tab:variables}
\centering
\footnotesize
\begin{tabular}{|p{2.2cm}|p{1.8cm}|p{2.4cm}|p{2.8cm}|}
\hline
\textbf{Variable} & \textbf{Tipo} & \textbf{Naturaleza} & \textbf{Indicador / Unidad} \\
\hline

Tiempo ($t$) & Independiente & Cuantitativa continua & Años \\
\hline

Actividad radiactiva ($A$) & Dependiente & Cuantitativa continua & Becquerel (Bq) \\
\hline

Concentración de Cs-137 & Dependiente & Cuantitativa continua & Bq/m$^2$ \\
\hline

Método de mitigación & Independiente & Cualitativa nominal & Tipo de intervención \\
\hline

Eficiencia de mitigación & Dependiente & Cuantitativa & Reducción (\%) \\
\hline

Tipo de suelo & Control & Cualitativa & Clasificación del suelo \\
\hline

Tipo de ecosistema & Control & Cualitativa & Clasificación ecológica \\
\hline

Vida media ($T_{1/2}$) & Parámetro & Constante física & Años \\
\hline

Constante de decaimiento ($\lambda$) & Parámetro & Constante física & año$^{-1}$ \\
\hline

\end{tabular}
\end{table}

\subsection{¿Bajo qué enfoques se analiza la realidad?}
El análisis se realiza bajo un enfoque mixto, combinando información cualitativa y cuantitativa, basado en fuentes especializadas como la IAEA sobre la gestión de accidentes nucleares severos.

\subsection{¿Quiénes son los sujetos o unidades de análisis?}
La población de estudio estará conformada por artículos científicos relacionados con la mitigación de la radiación y la respuesta ante desastres nucleares. Se seleccionará una muestra mínima de 15 artículos relevantes.

\subsection{¿Mediante qué herramientas se obtiene la evidencia?}
La evidencia se obtendrá mediante motores de búsqueda académicos, redes de citación (grafos), bibliotecas digitales y revistas científicas especializadas.

\subsection{¿Por qué este diseño es el más idóneo?}
Este diseño permite recopilar información de múltiples fuentes y autores, facilitando el análisis comparativo. Además, integra perspectivas de distintos contextos geográficos y temporales, lo que enriquece la investigación y mejora la comprensión del fenómeno estudiado.
